\section{Prujinali mayatnik}

\begin{enumerate}
    \item Bikrligi 10 N/m bo‘lgan prujinaga qanday massali (kg) yuk osilsa, tebranish davri 1 s bo‘lgan mayatnik hosil bo‘ladi? $\pi^2=10$.
    \begin{itemize}
        \item[A)] 10.
        \item[B)] 0,25.
        \item[C)] 2,5.
        \item[D)] 0,025.
        
    \end{itemize}

    \item Prujinaga osilgan 200 g massali yuk bir minutda 30 marta tebranadi. Prujinaning bikrligi necha N/m? $\pi^2=10$. [1]
    \begin{itemize}
        \item[A)] 0,2.
        \item[B)] 0,05.
        \item[C)] 0,3.
        \item[D)] 0,63.
        
    \end{itemize}

    \item Bikrligi 200 N/m bo‘lgan prujinaga osilgan yuk 8 s da 10 marta tebrandi. Yukning massasi qanday (kg)? $\pi^2=10$.
    \begin{itemize}
        \item[A)] 3,6.
        \item[B)] 3,2.
        \item[C)] 2,5.
        \item[D)] 1,6.
        
    \end{itemize}

    \item Prujinali mayatnik, og‘irlik kuchi tezlanishi Yernikidan 4 marta katta bo‘lgan planetaga Yerdan olib o‘tildi. Bunda uning tebranish davri qanday o‘zgaradi?
    \begin{itemize}
        \item[A)] o‘zgarmaydi.
        \item[B)] 2 marta ortadi.
        \item[C)] 4 marta ortadi.
        \item[D)] 2 marta kamayadi.
        
    \end{itemize}

    \item Agar tebranishlar davri $T$ bo‘lgan prujinali mayatnikka osilgan yukning massasi $n$ marta oshirilsa, tebranishlar davri qanday bo‘ladi?
    \begin{itemize}
        \item[A)] $nT$
        \item[B)] $n^2T$
        \item[C)] $\frac{T}{\sqrt{n}}$
        \item[D)] $\sqrt{n}T$
        
    \end{itemize}

    \item Prujinali tebrangich (mayatnik) dagi yukning massasi 9 marta ortsa, bo‘ylama tebranishlar davri qanday o‘zgaradi?
    \begin{itemize}
        \item[A)] 3 marta kamayadi.
        \item[B)] 3 marta ortadi.
        \item[C)] o‘zgarmaydi.
        \item[D)] 9 marta ortadi.
        
    \end{itemize}

    \item Agar prujina 2 marta qisqartirilib, yukning massasi 2 marta orttirilsa, prujinali mayatnikning tebranish davri qanday o‘zgaradi? [3]
    \begin{itemize}
        \item[A)] o‘zgarmaydi.
        \item[B)] 2 marta kamayadi.
        \item[C)] 2 marta ortadi.
        \item[D)] 4 marta kamayadi.
        
    \end{itemize}

    \item 1-rasmdagi sistemaning tebranish davri 10 s, 2-rasmda esa 8 s. 3-rasmdagi yukning tebranish davrini toping (s). Barcha prujinalarning bikirligi bir xil.

\begin{tikzpicture}
    % Umumiy uslublar
    \tikzset{
        wall/.style={fill, pattern=north east lines, draw=none},
        spring/.style={thick, decoration={coil, aspect=0.4, segment length=5pt, amplitude=3pt}, decorate},
        mass/.style={thick, fill=gray!30}
    }

    % --- 1-RASM (m1 va m2 vertikal) ---
    % Shift
    \draw[wall] (-0.8,0) rectangle (0.8, 0.2);
    \draw[thick] (-0.8,0) -- (0.8,0);
    % Prujina
    \draw[spring] (0, 0) -- (0, -2);
    % m1 yuk
    \draw[mass] (-0.3, -2) rectangle (0.3, -2.6);
    \node at (0.6, -2.3) {$m_1$};
    % Bog'lovchi ip
    \draw[thick] (0, -2.6) -- (0, -2.8);
    % m2 yuk
    \draw[mass] (-0.3, -2.8) rectangle (0.3, -3.4);
    \node at (0.6, -3.1) {$m_2$};
    \node[below] at (0, -3.5) {1-rasm};

    % --- 2-RASM (faqat m1) ---
    \begin{scope}[xshift=3cm]
        % Shift
        \draw[wall] (-0.8,0) rectangle (0.8, 0.2);
        \draw[thick] (-0.8,0) -- (0.8,0);
        % Prujina
        \draw[spring] (0, 0) -- (0, -2);
        % m1 yuk
        \draw[mass] (-0.3, -2) rectangle (0.3, -2.6);
        \node at (0.6, -2.3) {$m_1$};
        \node[below] at (0, -2.7) {2-rasm};
    \end{scope}

    % --- 3-RASM (m2 gorizontal) ---
    \begin{scope}[xshift=6cm, yshift=-1.5cm]
        % Devor (chapda)
        \draw[wall] (-0.2,-0.2) rectangle (0, 0.8);
        \draw[thick] (0,-0.2) -- (0,0.8);
        % Sirt (pastda)
        \draw[thick] (0,-0.2) -- (2.5,-0.2);
        % Prujina
        \draw[spring] (0, 0.15) -- (1.5, 0.15);
        % m2 yuk
        \draw[mass] (1.5, -0.2) rectangle (2.1, 0.4);
        \node at (1.8, 0.6) {$m_2$};
        \node[below] at (1.1, -0.3) {3-rasm};
    \end{scope}

\end{tikzpicture}
    
    \begin{itemize}
        \item[A)] 4.
        \item[B)] 5.
        \item[C)] 6.
        \item[D)] 9.
        
    \end{itemize}

    \item Rasmda tasvirlangan prujinali mayatniklarning tebranish davrlari nisbati $T_1/T_2$ ni toping?

\begin{tikzpicture}
    % Umumiy uslublar
    \tikzset{
        wall/.style={fill, pattern=north east lines, draw=none},
        spring/.style={thick, decoration={coil, aspect=0.4, segment length=4pt, amplitude=3pt}, decorate},
        mass/.style={thick, fill=gray!30}
    }

    % --- 1-SISTEMA (Parallel ulangan) ---
    \node at (0, 0.5) {\textbf{1)}};
    % Shift (shiftok)
    \draw[wall] (-1,0) rectangle (1, 0.2);
    \draw[thick] (-1,0) -- (1,0);

    % Prujinalar
    \draw[spring] (-0.6, 0) -- (-0.6, -2);
    \node[left] at (-0.7, -1) {$k$};
    
    \draw[spring] (0.6, 0) -- (0.6, -2);
    \node[right] at (0.7, -1) {$3k$};

    % Bog'lovchi gorizontal chiziq va yuk
    \draw[thick] (-0.6, -2) -- (0.6, -2);
    \draw[thick] (0, -2) -- (0, -2.3);
    \draw[mass] (-0.3, -2.3) rectangle (0.3, -2.9);
    \node at (0, -2.6) {$m$};

    % --- 2-SISTEMA (Ketma-ket ulangan) ---
    \begin{scope}[xshift=4cm]
        \node at (0, 0.5) {\textbf{2)}};
        % Shift (shiftok)
        \draw[wall] (-1,0) rectangle (1, 0.2);
        \draw[thick] (-1,0) -- (1,0);

        % Birinchi prujina
        \draw[spring] (0, 0) -- (0, -1.2);
        \node[right] at (0.1, -0.6) {$2k$};

        % Ikkinchi prujina
        \draw[spring] (0, -1.2) -- (0, -2.4);
        \node[right] at (0.1, -1.8) {$2k$};

        % Yuk
        \draw[mass] (-0.4, -2.4) rectangle (0.4, -3.1);
        \node at (0, -2.75) {$4m$};
    \end{scope}

\end{tikzpicture}
    
    % [Bu yerda rasm: 1) k va 3k parallel m bilan; 2) 2k va 2k ketma-ket 4m bilan]
    \begin{itemize}
        \item[A)] 4.
        \item[B)] 1/4.
        \item[C)] 2.
        \item[D)] 1/2.
        
    \end{itemize}

    \item Prujina o‘ziga osilgan yuk ta’sirida 2,5 cm ga cho‘zildi. Shu prujinali mayatnikning tebranish chastotasini aniqlang (Hz).
    \begin{itemize}
        \item[A)] 1,6.
        \item[B)] 25.
        \item[C)] 3,2.
        \item[D)] 5.
        
    \end{itemize}

    \item Quyidagi $m$ massali yukning tebranish davrini aniqlang. 

    \begin{tikzpicture}
    % Chap va o'ng devorlar uchun uslub
    \tikzset{wall/.style={fill, pattern=north east lines, draw=none}}
    
    % Prujina uslubi
    \tikzset{spring/.style={thick, decoration={coil, aspect=0.4, segment length=4pt, amplitude=3pt}, decorate}}

    % Chap devor
    \draw[wall] (0,-1) rectangle (-0.3, 1);
    \draw[thick] (0,-1) -- (0,1);

    % O'ng devor
    \draw[wall] (6,-1) rectangle (6.3, 1);
    \draw[thick] (6,-1) -- (6,1);
    \draw[thick] (0,-0.5) -- (6,-0.5);
    % Massali yuk (m)
    \draw[thick, fill=gray!20] (2.5,-0.5) rectangle (3.5, 0.5);
    \node at (3,0) {\large m};

    % Chap tarafdagi ikkita parallel prujina (k)
    \draw[spring] (0, 0.25) -- (2.5, 0.25);
    \node[above] at (1.25, 0.4) {$k$};
    
    \draw[spring] (0, -0.25) -- (2.5, -0.25);
    \node[above] at (1.25, -0.1) {$k$};

    % O'ng tarafdagi bitta prujina (k)
    \draw[spring] (3.5, 0) -- (6, 0);
    \node[above] at (4.75, 0.15) {$k$};

\end{tikzpicture}
    % [Bu yerda rasm: chapda ikkita k parallel, o'ngda bitta k]
    \begin{itemize}
        \item[A)] $T = 2\pi \sqrt{\frac{3m}{2k}}$
        \item[B)] $T = 2\pi \sqrt{\frac{2m}{3k}}$
        \item[C)] $T = 2\pi \sqrt{\frac{3m}{k}}$
        \item[D)] $T = 2\pi \sqrt{\frac{m}{3k}}$
        
    \end{itemize}
\end{enumerate}




\begin{center}
\begin{tabular}{|c|c|c|c|c|}
\hline
1 & 2 & 3 & 4 & 5 \\
\hline
B & A & C & A & D \\
\hline
6 & 7 & 8 & 9 & 10 \\
\hline
B & A & B & D & B \\
\hline
11 & & & & \\
\hline
D & & & & \\
\hline
\end{tabular}
\end{center}